% Generated draft: requires mathematical and source review.
\documentclass[11pt]{article}
\usepackage[margin=25mm]{geometry}
\usepackage{fontspec}
\setmainfont{Latin Modern Roman}
\usepackage{amsmath,amssymb,mathtools,unicode-math}
\setmathfont{Latin Modern Math}
\usepackage{enumitem,needspace,xurl,hyperref}
\setlength{\parindent}{0pt}
\setlength{\parskip}{5pt}
\setlength{\emergencystretch}{3em}
\newcommand{\N}{\mathbb{N}}
\newcommand{\Z}{\mathbb{Z}}
\newcommand{\Q}{\mathbb{Q}}
\newcommand{\R}{\mathbb{R}}
\newcommand{\C}{\mathbb{C}}
\newcommand{\F}{\mathbb{F}}
\newcommand{\A}{\mathbb{A}}
\newcommand{\PP}{\mathbb{P}}
\newcommand{\E}{\mathbb{E}}
\newcommand{\eps}{\varepsilon}
\newcommand{\dd}{\,\mathrm{d}}
\begin{document}
\section*{UnsolvedMath ID: 30004676}
\textbf{Source problem code:} catalogue entry
\subsection{Definitions and mathematical statement}
Let \(\operatorname{Adm}(A)\) mean that \(A\) is an admissible set in the sense intended by the source. Write \(D_\Sigma(B,A)\) for the assertion that the structure \(B\) is \(\Sigma\)-definable in \(A\). Retain the source notation \(J(A)\) for the jump structure of \(A\), \(\leq_\Sigma\) for its reducibility relation, and \(\equiv_\Sigma\) for its equivalence relation.

The supplied conjecture is
\[
\forall A\,\Bigl(\operatorname{Adm}(A)\Longrightarrow
\bigl[A\equiv_\Sigma J(A)
\ \Longleftrightarrow\ 
\forall B\,\bigl(D_\Sigma(B,A)\Longrightarrow B\leq_\Sigma A\bigr)\bigr]\Bigr),
\]
where \(B\) ranges over structures that are \(\Sigma\)-definable in \(A\).

\emph{Unresolved definitions and scope.} The supplied record does not define the admissibility convention, the class of structures and signatures, \(\Sigma\)-definability, \(\Sigma\)-reducibility, \(\Sigma\)-equivalence, or the construction \(J\). In particular, it does not specify whether parameters are permitted in definitions, how definable structures are represented, or how the admissible set \(A\) is viewed as a structure. These choices must be recovered from the cited source before this displayed formulation can be treated as a fully specified mathematical statement. No particular interpretation of those missing conventions is imposed here.
\subsection{Short English statement}
The conjecture asks whether, for every admissible set \(A\), being \(\Sigma\)-equivalent to its jump structure is equivalent to every structure \(\Sigma\)-definable in \(A\) being \(\Sigma\)-reducible to \(A\). The technical conventions are absent from the supplied record and require source review.
\subsection{Sources}
\begin{itemize}
\item Supplied local UnsolvedMath registry record, permanent ID \(30004676\), title ``Admissible Sets and Their Jump Structures,'' exact source problem code \texttt{catalogue entry}. The collection is ``Oberwolfach Reports Open Problems.'' No individual proposer is identified; the supplied proposed year is 2021.
\item Supplied report citation: \emph{Computability Theory (2022)}, Oberwolfach Reports, DOI 10.4171/owr/2021/21, pp.\ 1149--1190. \url{https://doi.org/10.4171/owr/2021/21}. The record describes Oberwolfach Reports as documenting mathematical workshops and their open-problem discussions.
\item Supplied literature source: ``Oberwolfach Report 21/2021.'' \url{https://ems.press/content/serial-article-files/46899}.
\item Supplied literature source: Puzarenko, \emph{Fixed points of the Sigma jump operator}. \url{https://old.math.nsc.ru/conference/malmeet/11/plenary/2011MM_Puzarenko.pdf}. This reference is supplied for context, without identifying Puzarenko as the proposer of this conjecture.
\item Literature assessment dated \texttt{2026-08-22}: the supplied curation labels the problem open and presents the equivalence as a conjecture. It states that fixed-point results for the jump operator provide context and examples, rather than establishing this equivalence for every admissible set. This is supplied dated curation derived from the OWR material, not an independently verified literature assessment; time-sensitive status and cited claims require independent verification.
\end{itemize}
\end{document}
